\section{Conclusion}
\label{sec:conclusion}

Epistract is the research-consolidation layer for scientists and computational analysts working inside Claude Code. Deep research happens upstream -- through Claude Code's built-in web tools, MCP-mediated source connectors, and the broader plugin ecosystem -- producing a body of documents that answers a research direction. Epistract takes that body and consolidates it into a two-layer knowledge graph, an epistemic classification of every relation, and an analyst briefing, from which the scientist answers a bundle of specific questions and iterates as their understanding of the consolidated body grows. Three architectural choices give the framework its current shape. First, a categorical \emph{epistemic super-domain layer} annotates every relation with a status (\texttt{asserted}, \texttt{prophetic}, \texttt{hypothesized}, \texttt{contested}, \texttt{contradiction}, \texttt{negative}, \texttt{speculative}) using a deterministic rule engine that combines hedging-pattern matching with document-type inference. Second, the framework is \emph{domain-agnostic}: a domain is four configuration files (schema, prompt, epistemic rules, persona), and the core pipeline does not need to be modified to add one. Third, the analyst persona is reified as a \emph{single source of truth} per domain that drives both a proactive narrator briefing and a reactive workbench chat, keeping the two surfaces coherent.

The framework ships with four pre-built domains. Three -- drug-discovery (six scenarios, 111 documents), clinicaltrials (one scenario, 10 protocols), and fda-product-labels (one scenario, 7 SPL labels) -- are validated. One -- contracts -- remains a schema scaffold awaiting its first showcase corpus. The framework documents this gap honestly and tracks it in the public issue queue rather than presenting all four domains as if at parity.

The framework's contribution is methodological as much as architectural. Treating the epistemic dimension as a first-class layer changes what a knowledge graph is for: not a structural artifact a human reads in a visualizer, but the input to an analyst persona that synthesizes briefings and answers reactive questions. The persona dual-use pattern -- one YAML string, two surfaces -- collapses what would otherwise be two separately-maintained system prompts into one and prevents drift. The controlled-port pattern (Section~\ref{sec:collaboration}) shows that the same architectural choices that make the framework pluggable also make external contribution tractable: a domain is a directory, not a patch against the core pipeline.

\paragraph{Scope: small-to-medium corpora that fit a Claude Code session.} The design point is interactive analytical sessions in which a scientist builds a bespoke graph to answer a specific question, iterates by adding documents and re-classifying, and follows up in chat. The seven validated scenarios range from 10 to 34 documents (Section~\ref{sec:evaluation}); the framework has not been tested at corpora of hundreds or thousands of documents and the upper bound at which the workbench chat or narrator briefing degrades is undetermined. The honest position is to scope the framework to interactive analytical work and to provide an explicit escape hatch for scale: \texttt{/epistract:export} dumps the graph to GraphML, CSV, SQLite, or JSON for ingestion into Neo4j-class persistent stores, where larger-scale or multi-team analysis can continue outside the Claude Code session.

\paragraph{Future directions.} Several directions remain open. Cross-scenario knowledge persistence (Issue~\#15) would let a domain compound across runs -- a canonical entity registry, then aggregate graphs, then learned rules -- and is the most architecturally meaningful next step. The fda-product-labels domain has had a first showcase corpus (S8) but the dominance of \texttt{reported}-tier classifications observed (97\%) suggests the four-level epistemology classifier is best exercised on a larger corpus that includes a higher proportion of CLINICAL STUDIES and CLINICAL PHARMACOLOGY sections; that is a candidate for an S8b or S9 expansion. \emph{MCP-mediated data-source connectors} are an aspirational direction worth highlighting: today a scientist assembles a corpus by ad-hoc query against PubMed, Google Scholar, ClinicalTrials.gov, or SerpAPI; tomorrow the same corpus could be assembled \emph{on demand} by Claude Code MCP servers exposing those sources (and proprietary sources within institutional environments) as structured tools available to the framework. This would close the loop between question and corpus: the scientist asks ``what does the corpus need to look like to answer this?'' and an MCP-enabled framework assembles a small corpus inside the session that fits the question. Additional domains beyond the four shipped will test whether the wizard's 15-minute claim holds across diverse corpora. And the integration with regulatory and clinical workflows -- the responsible-use boundary -- needs continuing attention as the framework's outputs move from research artifacts to decision inputs.

The framework, all corpora, all artifacts, and this paper are reproducible from \url{https://github.com/usathyan/epistract} under the MIT license.
